%%%%%%%%%%%%%%%%%%%%%%%%%%%%%%%%%%%%%%%%%%%%%%%%%%%%%%%%%%%%%%%%%%%%%%%%%%%%%%%%
% MedHive: A Proof-of-Concept Dashboard for Preventable Medication Waste
% Detection in Skilled Nursing Facilities
% Authors: Srinjoy Chatterjee, Shreya Dandwate, Timothy Fraser
% Cornell University -- Meinig School of Biomedical Engineering
%           & Department of Systems Engineering
%%%%%%%%%%%%%%%%%%%%%%%%%%%%%%%%%%%%%%%%%%%%%%%%%%%%%%%%%%%%%%%%%%%%%%%%%%%%%%%%

\documentclass[letterpaper, 10pt, conference]{IEEEconf}

\providecommand{\IEEEoverridecommandlockouts}{}
\providecommand{\overrideIEEEmargins}{}
\IEEEoverridecommandlockouts
\overrideIEEEmargins

\usepackage[utf8]{inputenc}
\usepackage[T1]{fontenc}
\usepackage{amsmath}
\usepackage{graphicx}
\usepackage{hyperref}
\usepackage{booktabs}
\usepackage{tabularx}
\usepackage{cite}
\usepackage{amssymb}

\title{\LARGE \bf
MedHive: A Proof-of-Concept Dashboard for Preventable Medication Waste
Detection in Skilled Nursing Facilities~---~A Customer Discovery and
Feasibility Study
}

\author{
Srinjoy Chatterjee$^{1}$, Shreya Dandwate$^{1}$, Timothy Fraser$^{1,2}$%
\thanks{$^{1}$Meinig School of Biomedical Engineering, Cornell University,
Ithaca, NY 14853, USA.
$^{2}$Department of Systems Engineering, Cornell University, Ithaca, NY 14853,
USA.
Emails: \texttt{sc2827@cornell.edu}, \texttt{sd2284@cornell.edu},
\texttt{tmf77@cornell.edu}}
}

\begin{document}

\maketitle
\thispagestyle{empty}
\pagestyle{empty}

%%%%%%%%%%%%%%%%%%%%%%%%%%%%%%%%%%%%%%%%%%%%%%%%%%%%%%%%%%%%%%%%%%%%%%%%%%%%%%%%
\begin{abstract}

\textbf{Objective:} Preventable prescription medication waste in skilled nursing
facilities (SNFs) represents a measurable and underaddressed cost burden on the
U.S.\ healthcare system, with annual losses in long-term care settings estimated
at \$2~billion or more.
This paper presents MedHive, a proof-of-concept R~Shiny dashboard designed to
visualize and flag preventable medication waste using Pareto analysis and
Statistical Process Control (SPC).
\textbf{Methods:} A structured customer discovery process was conducted across
six SNFs in the Ithaca, NY region via sequential outreach.
Three facilities participated in semi-structured informational interviews to
assess real-world data availability and validate synthetic data assumptions.
\textbf{Results:} None of the three interviewed facilities tracked
medication-level waste data in a standardized, exportable digital format.
Waste documentation was primarily confined to paper Medication Administration
Record (MAR) annotations and informal nurse notes.
Dashboard architecture was revised to relax schema requirements and support
imputed cost fields.
\textbf{Conclusion:} MedHive demonstrates technical feasibility as a
lightweight, open-access waste tracking tool; however, significant gaps between
modeled data assumptions and actual facility infrastructure must be addressed
before clinical deployment.
\textbf{Significance:} This work establishes a foundation for scalable
medication waste analytics in long-term care settings, with implications for
cost reduction, regulatory compliance, and prescriber feedback.

\end{abstract}

%%%%%%%%%%%%%%%%%%%%%%%%%%%%%%%%%%%%%%%%%%%%%%%%%%%%%%%%%%%%%%%%%%%%%%%%%%%%%%%%
\section{Introduction}

Prescription medication waste in skilled nursing facilities (SNFs) constitutes
a substantial and largely invisible financial burden on the U.S.\ healthcare
system.
Bazalo and Weiss estimated that unused solid oral medications dispensed to
Medicare Part~D residents in nursing facilities cost approximately \$125~million
annually, representing 3\% of total dispensed cost~\cite{Bazalo2011Measurement}.
More comprehensive estimates that account for all dosage forms and waste sources
place the total annual cost of unexpired drug waste in long-term care at
approximately \$2~billion, with broader system-wide losses---including hospital
and outpatient settings---approaching \$5~billion per year~\cite{Lenzer2014US}.
Early investigations of drug discard patterns in long-term care facilities
documented these losses as far back as 1977~\cite{Ciullo1977Drug}, yet
systematic, facility-level waste tracking tools remain largely absent from
clinical practice.

The root causes of SNF medication waste are well characterized in the
literature: prescriber discontinuation within days of initiating therapy,
patient death or transfer while on long-supply fills, non-adherence, and
administrative errors collectively account for the majority of
preventable events~\cite{Bekker2018Patient, Morgan2001The}.
Despite this knowledge, few of the approximately 15,000 certified SNFs in the
United States maintain structured, electronic records of waste at the
medication or reason-code level~\cite{Chambers2016Assessing}.
Electronic health record (EHR) adoption in long-term care lags far behind
that of acute care settings, and the facilities that have adopted EHR systems
often lack the data granularity or interoperability needed to export waste
analytics~\cite{Kruse2015Adoption, Vest2019Adoption}.

This paper presents MedHive, a free, open-access R~Shiny dashboard designed
to address this gap.
MedHive accepts user-uploaded CSV data conforming to a defined schema and
generates actionable waste insights through Pareto charts, SPC control charts,
rule-based operational alerts, and a first-fill policy simulator.
The system was developed under a Six Sigma DMAIC framework and validated
against three synthetic facility datasets representing high-waste, low-waste,
and balanced operational profiles.
To ground design decisions in operational reality, we conducted a structured
customer discovery process targeting SNFs in the Tompkins County, NY region.

The remainder of this paper is organized as follows.
Section~II reviews relevant literature on medication waste, EHR adoption in
SNFs, and healthcare applications of SPC and Six Sigma.
Section~III describes the customer discovery protocol and dashboard
architecture.
Section~IV presents findings from both the discovery interviews and an
illustrative case study using synthetic data.
Section~V discusses feasibility, limitations, and directions for future work,
and Section~VI concludes.

%%%%%%%%%%%%%%%%%%%%%%%%%%%%%%%%%%%%%%%%%%%%%%%%%%%%%%%%%%%%%%%%%%%%%%%%%%%%%%%%
\section{Literature Review}

\subsection{Medication Waste in Skilled Nursing Facilities}

The problem of prescription medication waste in long-term care is neither new
nor trivial.
Early observational studies documented drug discard patterns in long-term care
facilities as early as the 1970s, with Ciullo and Shepherd reporting that
central nervous system agents, anti-infectives, and cardiac drugs collectively
accounted for approximately 75\% of all Medicaid prescription
discards~\cite{Ciullo1977Drug}.
Brown and Kirk subsequently demonstrated that the introduction of unit-dose
drug distribution systems significantly reduced per-patient discard costs in
Indiana intermediate-care and SNFs~\cite{Brown1984Cost}, establishing an early
precedent for process-level intervention.

Despite these early insights, the scale of the problem grew substantially with
the aging of the U.S.\ population.
Morgan estimated that wasted prescription medication among outpatient adults
over 65 cost an average of \$30.47 per person per year---modest at the
individual level, but projecting to over \$1~billion annually for the national
cohort~\cite{Morgan2001The}.
Bazalo and Weiss subsequently quantified waste specifically within the Medicare
Part~D SNF population, finding that 6.8\% of all dispensed prescriptions were
returned partially used and that the cost of unconsumed solid oral medications
alone reached approximately \$125~million per year~\cite{Bazalo2011Measurement}.
Lenzer placed these figures in a broader context, reporting that long-term
care facilities waste roughly \$2~billion in unexpired medicines annually, with
system-wide estimates reaching \$5~billion when outpatient and hospital settings
are included~\cite{Lenzer2014US}.

Subsequent work has sought to characterize which patient and medication factors
drive preventable waste.
Bekker and colleagues found that over one-third of medications returned to Dutch
community pharmacies could be classified as preventable waste, with longer
dispensing periods---particularly supplies exceeding three months---carrying
significantly elevated risk~\cite{Bekker2018Patient}.
Jang and Kang confirmed this pattern in Korean long-term care data, finding
that 33.7\% of residents had unused medicines and that residents who died during
the study period were 15 times more likely to generate waste than
survivors~\cite{Jang2022Factors}.
Chambers and Austin highlighted a structural data gap underlying all of these
findings: few U.S.\ long-term care facilities maintain records of the
pharmaceutical waste they generate, making system-level estimates
uncertain~\cite{Chambers2016Assessing}.

Responses to this waste burden have ranged from redistribution protocols to
formal recycling programs.
Toh and Chew demonstrated that 90.8\% of donated medications from hospital and
clinic sources were reusable, yielding meaningful cost
savings~\cite{Toh2017Turning}, while Drummond and colleagues found that
recycling unused oral solid medications across a 21-hospital network could
divert approximately 461,000 units from incineration with an estimated net
value of \$415,000 per year~\cite{Drummond2023Recycling}.
Smale and colleagues argued that sustainable reduction requires joint
responsibility across manufacturers, prescribers, pharmacists, and health
authorities rather than any single downstream
intervention~\cite{Smale2021Waste-minimising}.
Critically, however, all of these downstream interventions presuppose the
existence of structured waste data at the facility level---a prerequisite
that remains largely unmet in the SNF setting.

\subsection{EHR Adoption and Data Infrastructure in Skilled Nursing Facilities}

Electronic health record adoption in long-term care settings has historically
lagged behind acute care, in large part because the Health Information
Technology for Economic and Clinical Health (HITECH) Act explicitly excluded
long-term care facilities from its financial
incentives~\cite{Kruse2015Adoption}.
Vest and colleagues found that even among facilities that had adopted health
information technology, interoperability was the most frequently cited barrier
to effective clinical information sharing, and only facilities classified as
``innovative'' showed adoption odds more than six times the
baseline~\cite{Vest2019Adoption}.
Felix and colleagues, examining all federally certified nursing homes in one
southern state, found that 69.8\% had some form of EHR in use, but that
adoption was concentrated among better-resourced facilities operating in
competitive markets---suggesting persistent inequities in digital
infrastructure~\cite{Felix2020Organizational}.

Systematic reviews have catalogued the barriers to EHR adoption in exhaustive
detail.
Kruse and colleagues identified 39 distinct barriers across 27 studies, with
cost, technical concerns, inadequate support structures, and resistance to
change appearing most
frequently~\cite{Kruse2016Barriers,Kruse2016Adoption}.
Ajami and Arab-Chadegani corroborated these findings, identifying resistance
to change as the single most significant implementation
obstacle~\cite{Ajami2013Barriers}.
Gesulga and colleagues further distinguished between people-centered barriers
(user resistance and skill deficits) and procedural barriers (concerns about
return on investment and insufficient administrative
support)~\cite{Gesulga2017Barriers}.
Tsai and colleagues synthesized 142 studies published between 2005 and 2020 and
found that EHR implementation produced both positive and negative effects on
clinical work, with poor training, resource constraints, and low technology
literacy consistently impeding full
adoption~\cite{Tsai2020Effects}.
Fennelly and colleagues, reviewing national EHR rollouts, identified 15
inter-linked organizational, human, and technological success factors,
emphasizing that contextual fit and end-user engagement were as important as
technical infrastructure~\cite{Fennelly2020Successfully}.

Where EHRs have been successfully deployed in SNFs, they have demonstrated
substantial analytical value.
Riester and colleagues used SNF EHR datasets to characterize patterns of
pro re nata analgesic prescribing and administration across 24,038 hip fracture
patients, demonstrating the feasibility of leveraging longitudinal EHR records
for pharmacoepidemiology in long-term care~\cite{Riester2024Use}.
Baughman and colleagues applied structured quality improvement methods to a
Veterans Affairs SNF and achieved measurable reductions in unintentional
medication discrepancies through comprehensive medication reconciliation
protocols~\cite{Baughman2021Improving}.
These successes, however, reflect best-case EHR environments; the majority of
SNFs---particularly smaller, rural, or under-resourced facilities---lack the
data infrastructure to replicate such analyses.

\subsection{Six Sigma and Statistical Process Control in Healthcare}

Statistical Process Control (SPC) was introduced to healthcare quality
improvement in a foundational 2003 paper by Benneyan, Lloyd, and Plsek, who
demonstrated that control charts provide a rigorous and accessible method for
separating common-cause variation from special-cause signals in clinical process
data~\cite{Benneyan2003Statistical}.
Thor and colleagues systematically reviewed 57 empirical studies published
between 1990 and 2004 and documented SPC applications across a wide range of
healthcare settings, cataloguing 12 categories of benefit and emphasizing
that the tool's power depends critically on correct application and contextual
risk adjustment~\cite{Thor2007Application}.

Subsequent work has extended both the theoretical foundations and the applied
scope of SPC in healthcare.
Fretheim and Tomic highlighted the underutilized potential of SPC data for
rigorous impact evaluation, arguing that time-series data routinely collected
in quality improvement programs could support interrupted time series analyses
if properly archived and published~\cite{Fretheim2015Statistical}.
Gupta and colleagues demonstrated the application of SPC in complex neonatal
intensive care environments, where non-standard data distributions and small
sample sizes pose particular methodological challenges~\cite{Gupta2023Challenging}.
Schmidtke and colleagues reported results from a cluster-randomized trial of
SPC training in English NHS hospitals, finding that chart use increased
substantially in both intervention and control arms---suggesting a broader
cultural diffusion of the methodology across the health
system~\cite{Schmidtke2024Cluster}.

The Six Sigma DMAIC framework---Define, Measure, Analyze, Improve,
Control---provides the organizational scaffold within which SPC tools are most
commonly deployed.
Monday traced the intellectual lineage of DMAIC from Shewhart's 1920s
statistical process theory through Deming's Plan-Do-Study-Act cycle to its
modern formalization at Motorola in the 1980s~\cite{Monday2022Define}.
Ahmed demonstrated that integrating DMAIC with the Theory of Constraints
can reduce medical costs, errors, and defects simultaneously by targeting
the binding bottleneck within a care process~\cite{Ahmed2019Integrating},
while Reis and colleagues found through systematic review that DMAIC
implementations in healthcare consistently reduced length of stay and improved
patient satisfaction~\cite{Reis2023DMAIC}.
Pharmacy applications have proven particularly receptive to Lean Six Sigma
methods: Sallam's review of 73 studies found that the most common outcomes
were decreased medication turnaround time, improved process efficiency, and
reduced medication errors~\cite{Sallam2024Enhancing}.
Thakur and colleagues confirmed these benefits in clinical laboratory settings,
where Lean and Six Sigma techniques reduced waste, variation, and work
imbalances~\cite{Thakur2022Lean}.
A systematic review by de~Mendon\c{c}a and colleagues of SPC applications in
abdominal surgery noted that while uptake is growing, methodological quality
varies and erroneous conclusions arise when the underlying analytical approach
is not rigorously applied~\cite{DeMendonca2024Use}---a caution applicable to
any new domain of deployment.

\subsection{Gap in the Literature}

Despite a well-documented epidemiology of medication waste in long-term care, a
growing body of SPC methodology, and three decades of EHR adoption research,
no lightweight, open-access analytics tool specifically designed for
medication-level waste tracking in SNFs currently exists.
Prior SPC applications have concentrated in acute care and in-hospital pharmacy
settings; EHR literature consistently identifies long-term care as the
least-resourced sector for digital data infrastructure; and medication waste
research documents the scale of the problem without providing operational
tooling for facility administrators.
MedHive is designed to occupy precisely this intersection.

%%%%%%%%%%%%%%%%%%%%%%%%%%%%%%%%%%%%%%%%%%%%%%%%%%%%%%%%%%%%%%%%%%%%%%%%%%%%%%%%
\section{Methods}

\subsection{Customer Discovery Protocol}

To ground MedHive's design in operational reality, a structured customer
discovery process was conducted targeting SNFs in the Tompkins County, NY
region during Fall~2025.
Six facilities were identified through CMS Nursing Home Compare and local
directories: Kendal at Ithaca, Beechtree Care Center, Cayuga Heights
Rehabilitation, the Reconstructive Home, Hospicare of Tompkins County, and Oak
Hill Manor.
Facilities were contacted via a sequential outreach funnel: (1)~cold telephone
call or voicemail, (2)~follow-up email with a one-paragraph project description,
and (3)~scheduling of a 15-minute semi-structured informational interview.

Each interview was conducted by telephone and targeted either the Director of
Nursing or the facility's pharmacy liaison.
A structured interview guide audited three primary assumptions embedded in
MedHive's initial synthetic dataset: (1)~whether facilities tracked
medication-level waste data digitally, (2)~what variables were routinely
captured (e.g., waste reason, drug cost, residual quantity), and
(3)~what format the data were stored in and whether export was feasible.
Interviews were recorded with verbal consent and analyzed thematically for
schema validation and design revision.

\subsection{Dashboard Architecture}

MedHive is implemented as a single-file R~Shiny application (\texttt{app.R})
using the packages \texttt{tidyverse}, \texttt{lubridate}, \texttt{scales},
\texttt{DT}, \texttt{janitor}, and \texttt{bslib}.
The data ingestion pipeline accepts one or more user-uploaded CSV files
conforming to the schema described in Table~\ref{table:schema}.
Each file is processed by the \texttt{read\_one\_csv()} function, which
applies column standardization, date parsing, and derivation of three
computed fields: per-unit cost (\texttt{unit\_cost}),
waste cost (\texttt{waste\_cost}), and reason class (\texttt{reason\_class}),
where \texttt{unit\_cost = cost\_usd / initial\_weight} and
\texttt{waste\_cost = waste\_units} $\times$ \texttt{unit\_cost}.

\begin{table}[h]
\caption{MedHive Minimum Required Data Schema}
\label{table:schema}
\begin{center}
\footnotesize
\begin{tabularx}{\columnwidth}{|p{2.7cm}|p{1.2cm}|X|}
\hline
\textbf{Variable} & \textbf{Type} & \textbf{Description} \\
\hline
\texttt{prescription\_id} & string & Unique prescription fill identifier \\
\hline
\texttt{patient\_id} & string & De-identified patient code \\
\hline
\texttt{medication} & string & Drug or drug class name \\
\hline
\texttt{scheduled\_date} & date & Intended dispense date (YYYY-MM-DD) \\
\hline
\texttt{refill\_date} & date & Actual dispense date (YYYY-MM-DD) \\
\hline
\texttt{cost\_usd} & numeric & Dispensed cost in USD \\
\hline
\texttt{prescriber} & string & Prescribing clinician name \\
\hline
\texttt{initial\_weight} & numeric & Total mass dispensed (grams) \\
\hline
\texttt{weight\_at\_stop} & numeric & Residual mass at therapy end (grams)$^{*}$ \\
\hline
\texttt{percent\_wasted} & numeric & Residual as \% of dispensed mass \\
\hline
\texttt{waste\_reason} & category & Controlled vocabulary (see text) \\
\hline
\end{tabularx}
\end{center}
\small{$^{*}$Optional following customer discovery; may be imputed.}
\end{table}

Waste reason codes are classified into three operational tiers by the
\texttt{reason\_class()} function: \textit{preventable} (medication change,
non-adherence, administrative error, lost medication, expired stock),
\textit{semi-preventable} (transfer to hospital), and
\textit{non-preventable} (deceased, no waste, other).
This taxonomy is used throughout the dashboard to filter views and scope
cost-saving estimates.

\subsection{Pareto Analysis}

For each filtered dataset, MedHive computes total waste cost aggregated by
both medication and prescriber dimensions.
Items are ranked in descending order of waste cost and a cumulative percentage
curve is overlaid, producing a dual-axis Pareto chart consistent with the
80/20 principle used throughout quality management.
The Pareto decomposition surfaces the small number of medications or
prescribers responsible for a disproportionate share of total waste cost,
enabling targeted intervention.

\subsection{SPC Control Chart}

MedHive aggregates waste cost to the weekly level and applies SPC control
limits via the \texttt{in\_control\_limits()} function.
Rather than the classical sample standard deviation, the implementation uses
a robust estimator based on the Median Absolute Deviation (MAD):

\begin{equation}
  \hat{\sigma} = 1.2533 \cdot \text{MAD}(x)
\label{eq:sigma}
\end{equation}

where $x$ is the vector of weekly waste costs and the scaling constant
$1.2533 \approx 1/\Phi^{-1}(3/4)$ recovers an asymptotically unbiased
estimate of the population standard deviation under normality while
remaining resistant to outliers~\cite{Benneyan2003Statistical}.
The center line, Upper Control Limit (UCL), and Lower Control Limit (LCL)
are then computed as:

\begin{equation}
  \text{UCL} = \bar{x} + 3\hat{\sigma}, \quad
  \text{LCL} = \bar{x} - 3\hat{\sigma}
\label{eq:ucllcl}
\end{equation}

Weekly observations exceeding the UCL are flagged as special-cause variation,
warranting investigation by facility staff.

\subsection{Operational Alerts}

The \texttt{flags\_tbl()} reactive module generates four categories of
rule-based alert.
(1)~\textit{Policy alerts} flag any medication for which preventable waste
constitutes more than 30\% of total waste cost and recommend trialing a
shorter first-fill supply.
(2)~\textit{Prescriber alerts} identify clinicians whose preventable waste
rate exceeds the facility peer median by more than 10~percentage points,
among those with at least 20 attributed prescriptions.
(3)~\textit{Workflow alerts} flag all waste events linked to patient death
where residual medication units are greater than zero, prompting review of
auto-halt and EHR notification workflows.
(4)~\textit{Transfer alerts} summarize waste attributable to hospital transfer
events and estimate a 20\% recoverable fraction through improved partial-fill
and handoff protocols.

\subsection{First-Fill Policy Simulator}

The \texttt{sim\_detail()} reactive models projected savings under a
partial-dispensing policy.
For each preventable waste event $i$, an initial supply of
$q'_i = \max\!\bigl(1,\,\lfloor w_{\text{init},i} \cdot f \rfloor\bigr)$
units is dispensed, where $f \in [0.25, 1.0]$ is the user-specified
first-fill fraction.
The reducible quantity and units saved are then:

\begin{align}
  q'_i &= \max\!\bigl(1,\,\lfloor w_{\text{init},i} \cdot f \rfloor\bigr)
  \label{eq:simulator} \\
  s_i  &= \min\!\bigl(w_{\text{waste},i},\;
           \max(0,\, w_{\text{init},i} - q'_i)\bigr) \notag
\end{align}

Total projected savings $S$ over the filtered period are:

\begin{equation}
  S = \sum_{i=1}^{n} s_i \cdot c_i \cdot
      \mathbb{1}\bigl[\text{reason\_class}_i = \textit{preventable}\bigr]
\label{eq:savings}
\end{equation}

where $c_i$ is the per-unit cost of prescription~$i$.
Results are disaggregated by medication and presented as an interactive
ranked bar chart alongside a downloadable detail table.

\subsection{Deployment}

MedHive is deployed as a stateless, zero-authentication application on
\texttt{shinyapps.io}, accessible at no cost to any SNF administrator with
a web browser.
The stateless architecture was selected deliberately for the proof-of-concept
phase: it eliminates the need for HIPAA-compliant session management,
credential storage, or server-side data persistence, allowing the team to
demonstrate core analytical functionality without incurring the security and
compliance overhead required for production clinical deployment.
The full source code and synthetic demonstration datasets are publicly available
at \url{https://github.com/chatterjee-srinjoy/Control-Freaks-Six-Sigma}.

%%%%%%%%%%%%%%%%%%%%%%%%%%%%%%%%%%%%%%%%%%%%%%%%%%%%%%%%%%%%%%%%%%%%%%%%%%%%%%%%
\section{Findings}

\subsection{Customer Discovery Results}

Of the six SNFs contacted, three agreed to participate in informational
interviews, yielding a response rate of 50\%.
The three non-responding facilities either did not return follow-up contact
or declined participation citing staff time constraints.
Interviews were completed with two Directors of Nursing and one pharmacy
liaison across three separate facilities.
Table~\ref{table:interviews} summarizes the facility profiles and key
thematic findings.

\begin{table}[h]
\caption{Customer Discovery Interview Summary (Tompkins County, NY, Fall 2025)}
\label{table:interviews}
\begin{center}
\footnotesize
\begin{tabularx}{\columnwidth}{|p{0.55cm}|p{2.4cm}|p{1.6cm}|X|}
\hline
\textbf{Fac.} & \textbf{EHR System} & \textbf{Role} & \textbf{Key Finding} \\
\hline
F1 & PointClickCare & Dir.\ of Nursing & Waste in paper MAR notes; no digital export \\
\hline
F2 & MatrixCare & Pharmacy Liaison & Returns tracked by count, not cost or reason code \\
\hline
F3 & Paper/Excel hybrid & Dir.\ of Nursing & Medication lists via weekly fax; no cost attribution \\
\hline
\end{tabularx}
\end{center}
\end{table}

Several consistent themes emerged across all three interviews.
First, none of the facilities tracked medication waste at the individual
prescription or reason-code level in any standardized digital format.
Waste events were documented informally in paper MAR annotations or nursing
shift notes, and this information was not routinely aggregated, costed, or
reviewed at an administrative level.
Second, all three facilities had some form of return-to-pharmacy workflow for
residual medications, but only one could estimate the volume of returned
items (by pill count), and none could report aggregate waste costs.
Third, the Director of Nursing at Facility~F1 identified adverse drug
reactions and prescriber-initiated medication changes as the most frequently
cited informal waste categories, consistent with the controlled vocabulary
used in MedHive's synthetic data.
Fourth, all three participants expressed interest in a free, browser-based
tool if it could accept data in a format their pharmacy system could generate---
though none were confident that their current systems could produce the schema
as originally specified.

\subsection{Data Availability Gap}

Table~\ref{table:gap} summarizes the discrepancy between variables assumed in
MedHive's original synthetic dataset and variables confirmed as available
through customer discovery interviews.

\begin{table}[h]
\caption{Synthetic Data Assumptions vs.\ Facility-Reported Availability}
\label{table:gap}
\begin{center}
\footnotesize
\begin{tabularx}{\columnwidth}{|p{2.6cm}|p{1.2cm}|X|}
\hline
\textbf{Variable} & \textbf{Assumed} & \textbf{Facility Reality} \\
\hline
\texttt{waste\_reason} & Yes & Informal text; no controlled vocabulary \\
\hline
\texttt{cost\_usd} & Yes & Not tracked per-prescription; billed to payer only \\
\hline
\texttt{weight\_at\_stop} & Yes & Not recorded; return tracked by pill count only \\
\hline
Digital format & Yes & Paper MAR (1 fac.); EHR non-exportable (2 fac.) \\
\hline
Prescriber link. & Yes & Available in EHR; not linked to waste events \\
\hline
\end{tabularx}
\end{center}
\end{table}

\subsection{Dashboard Iteration}

Based on interview findings, three adjustments were made to MedHive's data
ingestion pipeline and schema requirements.
First, \texttt{weight\_at\_stop} was reclassified as an optional field;
when absent, the dashboard estimates residual quantity from supply length and
standard dispensing unit data.
Second, \texttt{cost\_usd} was made imput\-able from public drug pricing
references (e.g., the Medispan Red Book) using medication name as a join key,
allowing facilities without per-prescription cost records to still generate
cost-denominated waste estimates.
Third, the \texttt{waste\_reason} controlled vocabulary was supplemented with
a free-text annotation field and a simplified three-tier classification
(preventable / semi-preventable / non-preventable) as a fallback for facilities
unable to provide reason codes.

\subsection{Illustrative Case Study}

To demonstrate MedHive's analytical outputs in a realistic operational
context, consider a skilled nursing facility comparable to Facility~A in the
synthetic dataset---a 75-bed SNF with high medication turnover due to frequent
transfers and therapeutic changes.
Applying MedHive to three months of synthetic data ($n = 302$ prescription
events, August--October 2025), the dashboard surfaces the following insights.

\begin{figure}[h]
\centering
\fbox{\parbox{0.95\linewidth}{\centering\vspace{1.5cm}
\textit{[FIGURE 1: Pareto chart --- medication by total waste cost.
Top three drugs account for $\approx$42\% of preventable waste.
Insert dashboard screenshot here.]}
\vspace{1.5cm}}}
\caption{Pareto analysis of medication waste cost, Facility~A synthetic
dataset ($n=302$, Aug--Oct 2025). The 80\% cumulative threshold is indicated
by the dashed reference line.}
\label{fig:pareto}
\end{figure}

Pareto analysis (Fig.~\ref{fig:pareto}) identifies Risperidone, Warfarin, and
Gabapentin as the top three contributors to preventable waste cost,
collectively accounting for approximately 42\% of total preventable expenditure.
All three are high-cost, high-residual medications frequently prescribed on
long supply durations---consistent with Bekker et al.'s finding that longer
dispensing periods significantly elevate preventable waste
risk~\cite{Bekker2018Patient}.

\begin{figure}[h]
\centering
\fbox{\parbox{0.95\linewidth}{\centering\vspace{1.5cm}
\textit{[FIGURE 2: SPC time-series chart of weekly waste cost with 3$\sigma$
control limits. Two out-of-control weeks visible. Insert dashboard screenshot.]}
\vspace{1.5cm}}}
\caption{SPC control chart of weekly waste cost, Facility~A. The dashed
center line represents $\bar{x}$; dotted red lines indicate the UCL and LCL
at $\bar{x} \pm 3\hat{\sigma}$ using the MAD-based robust estimator
(Eq.~\ref{eq:sigma}).}
\label{fig:spc}
\end{figure}

The SPC control chart (Fig.~\ref{fig:spc}) reveals two weeks in which weekly
waste cost exceeded the upper control limit, consistent with special-cause
variation attributable to a cluster of same-week medication change orders.
The use of the MAD-based $\hat{\sigma}$ estimator (Eq.~\ref{eq:sigma}) ensures
that these UCL/LCL bounds are not artificially widened by the very outlier
weeks they are intended to detect.

\begin{figure}[h]
\centering
\fbox{\parbox{0.95\linewidth}{\centering\vspace{1.5cm}
\textit{[FIGURE 3: First-fill policy simulator output at $f = 0.5$.
Estimated savings by medication. Insert dashboard screenshot.]}
\vspace{1.5cm}}}
\caption{Projected first-fill policy savings by medication, Facility~A,
at first-fill fraction $f = 0.50$. Only preventable waste events are
included in the savings calculation (Eq.~\ref{eq:savings}).}
\label{fig:simulator}
\end{figure}

Applying the first-fill simulator at $f = 0.50$ (Fig.~\ref{fig:simulator})
projects a 50\% partial dispensing policy that, applied to preventable events
only, would reduce waste cost by an estimated 28--34\% for the highest-impact
medications.
The operational alert module additionally flags two prescribers whose
preventable waste rates exceed the facility peer median by more than
10~percentage points, providing a starting point for targeted prescriber
education and feedback.

%%%%%%%%%%%%%%%%%%%%%%%%%%%%%%%%%%%%%%%%%%%%%%%%%%%%%%%%%%%%%%%%%%%%%%%%%%%%%%%%
\section{Discussion}

\subsection{Feasibility Assessment}

MedHive's technical architecture is sound and its analytical outputs are
grounded in established quality improvement methodology.
The Pareto and SPC tools are standard, well-validated components of the Six
Sigma DMAIC framework~\cite{Benneyan2003Statistical, Thor2007Application,
Sallam2024Enhancing}, and the first-fill policy simulator provides a
transparent, auditable model for projecting the cost impact of partial-dispensing
policies prior to implementation.
The dashboard's stateless, browser-based deployment model minimizes barriers
to adoption for facilities that lack dedicated IT infrastructure---a deliberate
design choice consistent with the EHR adoption literature's emphasis on cost
and implementation complexity as primary
barriers~\cite{Kruse2016Barriers,Tsai2020Effects}.

The central feasibility challenge, however, is not technical but
infrastructural: the tool is only as useful as the data that can be supplied
to it.
Customer discovery revealed that none of the three interviewed facilities
possessed the data granularity assumed in MedHive's original schema.
This finding is not surprising in light of the EHR adoption
literature~\cite{Kruse2015Adoption, Vest2019Adoption, Felix2020Organizational},
which consistently documents that long-term care facilities are the
least-resourced sector for structured health data.
What the interviews add to this picture is specificity: even facilities using
widely adopted commercial SNF EHR platforms (PointClickCare, MatrixCare) do
not configure those systems to capture waste reason codes or per-prescription
costs at a level exportable to a third-party analytics tool.

The schema relaxations implemented following customer discovery---making
\texttt{weight\_at\_stop} optional and enabling cost imputation from
public pricing databases---meaningfully lower the barrier to entry.
A facility with only a medication name, prescriber, and dispense date can now
generate partial analytics; a facility with return-pill-count data can generate
Pareto rankings even without unit cost.
The minimum viable dataset for meaningful SPC output is weekly prescription
volume by medication---a field available from the medication administration
records of all three interviewed facilities.

\subsection{Limitations}

This study has several important limitations.
First, customer discovery was limited to three facilities in a single small
metropolitan area.
Tompkins County, NY is home to a research university and may not be
representative of the national SNF landscape in terms of administrator
sophistication, data awareness, or technological infrastructure.
Findings should be interpreted as directional rather than
generalizable~\cite{Chambers2016Assessing}.
Second, all performance claims in this paper are based on synthetic data.
The three facility datasets were constructed to represent plausible
operational profiles (high-waste, low-waste, and balanced), but they have
not been validated against actual facility records.
Waste cost estimates produced by the policy simulator carry no clinical or
financial guarantee.
Third, MedHive has not undergone usability testing with SNF staff.
The interface was designed by engineers and has not been evaluated by the
Directors of Nursing, pharmacy liaisons, or administrators who would be its
primary users.
Fourth, as a stateless application, MedHive currently stores no data between
sessions, preventing longitudinal trend analysis or multi-period benchmarking
without re-upload.
Fifth, the operational alert thresholds---30\% preventable share for policy
alerts, 10~percentage points above peer median for prescriber alerts---were
selected by the development team based on judgment and the structure of the
synthetic data.
These parameters have not been empirically validated and should be treated as
configurable defaults.

\subsection{Future Work}

Four directions for future development are prioritized.
First, direct EHR API integration with major SNF platforms (PointClickCare,
MatrixCare, Yardi Senior Living) would eliminate the CSV upload requirement
and allow real-time dashboarding without manual data export~\cite{Riester2024Use}.
Second, HIPAA-compliant authentication and encrypted session management are
prerequisites for any deployment involving real patient data; a Supabase or
Firebase backend with role-based access control would address this
requirement while preserving the tool's low-cost deployment model.
Third, multi-facility benchmarking---allowing a facility to compare its waste
metrics against an anonymized peer cohort---would substantially increase the
tool's operational value and is a natural extension of the existing
multi-file upload capability.
Fourth, and most critically, prospective clinical validation in partnership
with at least one SNF operator is necessary to establish that MedHive's
outputs are interpretable, actionable, and accurate when applied to real
facility data.
Such a study would simultaneously validate the revised schema, test the
operational alert thresholds, and provide the first empirical cost-impact
estimates outside of simulation.

%%%%%%%%%%%%%%%%%%%%%%%%%%%%%%%%%%%%%%%%%%%%%%%%%%%%%%%%%%%%%%%%%%%%%%%%%%%%%%%%
\section{Conclusions}

This paper has presented MedHive, a proof-of-concept R~Shiny dashboard for
preventable medication waste detection in skilled nursing facilities.
Developed under a Six Sigma DMAIC framework, MedHive integrates Pareto
analysis, MAD-based SPC control charts, rule-based operational alerting, and a
first-fill policy simulator into a unified, zero-cost, browser-accessible
interface.

A structured customer discovery process conducted across six SNFs in Tompkins
County, NY found that none of the three participating facilities tracked
medication waste at the individual prescription or reason-code level in an
exportable digital format---confirming that the primary barrier to adoption
is not tool availability but data infrastructure.
In response, the dashboard schema was revised to support optional and imputed
fields, substantially lowering the minimum data requirements for useful output.

Illustrative case study analysis using synthetic facility data demonstrated
that MedHive can surface actionable insights: Pareto decomposition identifies
the small number of medications driving disproportionate waste cost;
SPC control charts detect anomalous waste weeks not attributable to routine
variation; and the policy simulator projects meaningful savings from even
modest first-fill fraction reductions.

MedHive is freely available at
\url{https://github.com/chatterjee-srinjoy/Control-Freaks-Six-Sigma}
as an open-access tool for the SNF community.
The authors hope that this proof-of-concept will stimulate further investment
in medication waste analytics infrastructure for long-term care, and that
future work will deliver a clinically validated, HIPAA-compliant version of
the dashboard in partnership with operating facilities.

% Uncomment and tune only on final draft to balance last page columns:
% \addtolength{\textheight}{-12cm}

%%%%%%%%%%%%%%%%%%%%%%%%%%%%%%%%%%%%%%%%%%%%%%%%%%%%%%%%%%%%%%%%%%%%%%%%%%%%%%%%
\section*{Acknowledgment}

The authors thank Dr.\ Timothy Fraser and the Cornell Department of Systems
Engineering for guidance throughout this project.
This work was conducted as part of the Cornell Meinig School of Biomedical
Engineering M.Eng.\ program.
The authors also thank the nursing home administrators and Directors of
Nursing in Tompkins County, NY who generously participated in informational
interviews.

%%%%%%%%%%%%%%%%%%%%%%%%%%%%%%%%%%%%%%%%%%%%%%%%%%%%%%%%%%%%%%%%%%%%%%%%%%%%%%%%
\bibliographystyle{IEEEtran}
\bibliography{references}

%%%%%%%%%%%%%%%%%%%%%%%%%%%%%%%%%%%%%%%%%%%%%%%%%%%%%%%%%%%%%%%%%%%%%%%%%%%%%%%%
\appendix
\section{Appendix: Minimum CSV Schema and Source Code}

The MedHive source code, synthetic demonstration datasets (Facilities~A, B,
and C), and the codebook documenting all schema variables are available at:
\begin{center}
\url{https://github.com/chatterjee-srinjoy/Control-Freaks-Six-Sigma}
\end{center}

The minimum required CSV fields for SPC output are:
\texttt{prescription\_id}, \texttt{medication}, \texttt{scheduled\_date},
and \texttt{waste\_reason}.
Full Pareto and cost-denominated outputs additionally require
\texttt{cost\_usd} (or an imputed cost via medication name lookup) and
\texttt{initial\_weight}.
The policy simulator requires \texttt{initial\_weight},
\texttt{weight\_at\_stop} (or an estimate), and \texttt{cost\_usd}.

\end{document}
